% !TeX root = ../tesis.tex
%%%%%%%%%%%%%%%%%%%%%%%%%%%%%%%%%%%%%%%%%%%%%%%%%%%%%%%%%%%%%%%%%%%
% SECCIÓN 5.1 — LA RED METROPOLITANA COMO OBJETO DE ANÁLISIS
%%%%%%%%%%%%%%%%%%%%%%%%%%%%%%%%%%%%%%%%%%%%%%%%%%%%%%%%%%%%%%%%%%%
\section{La Red Metropolitana como Objeto de Análisis}
\label{sec:intro-resultados}

La construcción del grafo topológico de la \zmvm constituye el punto de partida
de todos los indicadores que se presentan en este capítulo. Para transformar la
geografía estática de las rutas de transporte en un modelo operativo y ruteable,
el \textbf{Modelo VFT} implementa una arquitectura de procesamiento en tres
capas progresivas, cuyo flujo se sintetiza en la Figura~\ref{fig:arq-capas}.

% --- Figura: arquitectura en 3 capas ---
\begin{figure}[H]
    \centering
    \begin{subfigure}[b]{0.32\textwidth}
        \centering
        \includegraphics[width=\textwidth]{Figures/Cap5/arq_capa1.png}
        \caption{Capa 1: Adquisición y procesamiento geométrico.}
        \label{fig:arq-capa1}
    \end{subfigure}
    \hfill
    \begin{subfigure}[b]{0.32\textwidth}
        \centering
        \includegraphics[width=\textwidth]{Figures/Cap5/arq_capa2.png}
        \caption{Capa 2: Modelado topológico y dinámico.}
        \label{fig:arq-capa2}
    \end{subfigure}
    \hfill
    \begin{subfigure}[b]{0.32\textwidth}
        \centering
        \includegraphics[width=\textwidth]{Figures/Cap5/arq_capa3.png}
        \caption{Capa 3: Evaluación y salida de indicadores.}
        \label{fig:arq-capa3}
    \end{subfigure}
    \caption[Arquitectura en tres capas del Modelo VFT]{Representación esquemática
    de las tres capas de procesamiento del Modelo VFT: adquisición geométrica,
    modelado topológico con impedancia temporal, y evaluación de indicadores.}
    \label{fig:arq-capas}
    \fuentefigura{Elaboración propia.}
\end{figure}

% --- Estadísticas del grafo ---
\subsection{Caracterización del Grafo Construido}
\label{subsec:estadisticas-grafo}

La ingesta de datos desde \texttt{Apimetro} proporcionó \textbf{22{,}766 entidades
espaciales} que, tras el proceso de \termino{Snapping Lógico Geográfico} mediante
un árbol \texttt{KDTree} con tolerancia de 50~m, se consolidaron en un grafo
dirigido con las siguientes dimensiones:

\begin{itemize}
    \item \textbf{11{,}115 nodos topológicos} (estaciones operativas, eliminando
    nodos fantasma y redundancias geométricas).
    \item \textbf{45{,}679 aristas}, de las cuales 25{,}197 corresponden a
    segmentos de servicio real y 20{,}482 a aristas de transbordo peatonal
    generadas algorítmicamente.
    \item \textbf{11 sistemas de transporte} representados: \textsc{Metro},
    Metrobús, Trolebús, Tren Ligero, Tren Suburbano, Cablebús, Mexicable,
    Mexibús, \textsc{rtp}, Corredores Concesionados y Pumabús.
    \item Distribución modal: \textbf{12\% infraestructura masiva estructurada}
    (Metro, MB, TL, CBB, SUB) y \textbf{88\% transporte de superficie}
    (RTP, CC, Trolebús, Mexibús, Mexicable).
\end{itemize}

La Figura~\ref{fig:paneles-grafo} ilustra las cuatro perspectivas de
visualización del grafo resultante: la infraestructura agregada, la conectividad
interestación, y la direccionalidad de ida y regreso.

% --- Figura: 4 paneles del grafo ---
\begin{figure}[H]
    \centering
    \begin{subfigure}[b]{0.48\textwidth}
        \centering
        \includegraphics[width=\textwidth]{Figures/Cap5/panel1_infraestructura.png}
        \caption{Panel 1: Infraestructura agregada de la red.}
        \label{fig:panel1}
    \end{subfigure}
    \hfill
    \begin{subfigure}[b]{0.48\textwidth}
        \centering
        \includegraphics[width=\textwidth]{Figures/Cap5/panel2_conectividad.png}
        \caption{Panel 2: Conectividad interestación (aristas).}
        \label{fig:panel2}
    \end{subfigure}

    \vspace{0.5em}

    \begin{subfigure}[b]{0.48\textwidth}
        \centering
        \includegraphics[width=\textwidth]{Figures/Cap5/panel3_ida.png}
        \caption{Panel 3: Direccionalidad de servicio (sentido ida).}
        \label{fig:panel3}
    \end{subfigure}
    \hfill
    \begin{subfigure}[b]{0.48\textwidth}
        \centering
        \includegraphics[width=\textwidth]{Figures/Cap5/panel4_regreso.png}
        \caption{Panel 4: Direccionalidad de servicio (sentido regreso).}
        \label{fig:panel4}
    \end{subfigure}
    \caption[Visualización del grafo topológico de la \zmvm]{Cuatro perspectivas
    de visualización del grafo dirigido construido sobre la red multimodal de la
    \zmvm. El proceso de Snapping Lógico redujo los nodos en aproximadamente
    un 90\% respecto a la representación geométrica cruda.}
    \label{fig:paneles-grafo}
    \fuentefigura{Elaboración propia con \texttt{VFTModel} (corrida 2026-04-18).}
\end{figure}

% --- Tabla resumen de indicadores ---
\subsection{Estructura de Indicadores y Dependencias por Fase}
\label{subsec:tabla-indicadores}

Los ocho indicadores que componen el análisis topológico se organizan en tres
fases cuya ejecución es secuencialmente dependiente, como se muestra en el
Cuadro~\ref{tab:resumen-indicadores-cap5}.

\begin{table}[H]
    \centering
    \caption[Resumen de indicadores del Modelo VFT por fase]{Resumen de los
    indicadores del Modelo VFT, su fase de cálculo, dependencias y estado
    de avance en la investigación.}
    \label{tab:resumen-indicadores-cap5}
    \begin{tabularx}{\textwidth}{|l|l|X|l|}
        \hline
        \textbf{Indicador} & \textbf{Fase} & \textbf{Dependencias} & \textbf{Estado} \\
        \hline
        Validación de Impedancia    & 1 & Grafo base + fallbacks            & Calculado \\
        Cobertura Espacial ($C$)    & 1 & Geometría de estaciones           & Calculado \\
        Fuerza Capilar ($FC$)       & 1 & Grafo base                        & Calculado \\
        Factor de Desviación ($DF$) & 1 & Grafo base + impedancia           & Calculado \\
        \hline
        Coef. de Fricción Vial ($CF$)  & 2 & Fase 1 completa                & Calculado \\
        Penalización por Transf. ($T_b$) & 2 & Fase 1 + $CF$              & Formulado \\
        Node Score ($VFT\_score$)   & 2 & Fase 1 + datos MV1/MV2            & Planificado \\
        \hline
        Tiempo de Viaje Prom. ($T$) & 3 & Fases 1 y 2 completas             & Pendiente \\
        Centralidad Interm. ($B$)   & 3 & Fases 1 y 2 + $T$                 & Pendiente \\
        Vulnerabilidad ($\Delta E$) & 3 & Fases 1, 2 y $B$                  & Pendiente \\
        \hline
    \end{tabularx}
    \fuentefigura{Elaboración propia.}
\end{table}
